\lecture{22}
\title{One-Parameter Subgroups, Lie Algebras, Lie Brackets}

\begin{document}

The previous few lectures have prepared us for the following definitions.

\textbf{Definition (Lie Group).} A \textbf{Lie group} is a group that is also a smooth manifold (i.e. looks like $\mathbb{R}^n$ locally).

\textbf{Example.} All of the following are Lie groups: $\{GL_n(\mathbb{R}), GL_n(\mathbb{C}), SO_n, U_n, SU_n\}$.  However, $GL_n(\mathbb{Q})$ is not a Lie group---it has too many ``holes'' in it.

\textbf{Definition (One-Parameter Subgroup).} Given a Lie group $G$, a \textbf{one-parameter subgroup} $H \subseteq G$ is a subgroup that is the image of a \textit{differentiable} homomorphism $\varphi: \mathbb{R} \to G$.

\textbf{Example.} The longitudinal subgroups from last lecture are one-parameter subgroups of $SU_2$.

\subsection{Matrix Exponentiation: One-Parameter $GL_n(\mathbb{C})$ Subgroups}

\textbf{Definition (Matrix Exponentiation).} Given any $A \in \mathbb{C}^{n \times n}$, we say $e^A = I_n + A + \frac{A^2}{2!} + \frac{A^3}{3!} + \dots$ (Taylor Series).

\begin{theorem}{Exponent Laws}
    For any $A, B \in \mathbb{C}^{n \times n}$ such that $AB = BA$, we have $e^{A + B} = e^A e^B$.
\end{theorem}

\begin{proof}
    Apply the distributive law to $e^Ae^B = \left(\sum_{k \geq 0} \frac{A^k}{k!}\right)\left(\sum_{\ell \geq 0} \frac{B^{\ell}}{\ell!}\right)$, where $AB = BA$ lets us collate factors together.

    Really, this theorem is just restating $(x + y)^n = \sum_{k = 0}^n \binom{n}{k}x^ky^{n - k}$, which only holds when $x$ and $y$ commute. ~ $\blacksquare$
\end{proof}

To actually compute the exponentiation by a matrix, use diagonalizability.

\begin{theorem}{Diagonalizable Computations}
    If $A$ is diagonalizable, then we have:
    \[ A = B \begin{bsmallmatrix} \lambda_1 & \cdots & 0 \\ \vdots & \ddots & \vdots \\ 0 & \cdots & \lambda_n \end{bsmallmatrix}B^{-1} ~~~ \implies ~~~ e^A = B \begin{bsmallmatrix} e^{\lambda_1} & \cdots & 0 \\ \vdots & \ddots & \vdots \\ 0 & \cdots & e^{\lambda_n} \end{bsmallmatrix} B^{-1}. \]  \vspace{-0.5cm}
\end{theorem}

In fact, \emph{most} matrices are diagonalizable, so\dots

\begin{theorem}{Determinant and Trace}
    For any $A \in \mathbb{C}^{n \times n}$, we have $\det e^A = e^{\mathrm{tr}(A)}$.
\end{theorem}

\begin{proof}
    This holds when $A$ is diagonalizable.  By continuity, it holds for all $A$. ~ $\blacksquare$
\end{proof}

The point of all this setup is the following observation:

\begin{boxed-text}[28em]
    \textbf{Observation.} For any $A \in \mathbb{C}^{n \times n}$, the map $\varphi_A: \mathbb{R} \to GL_n(\mathbb{C})$ via $\varphi_A(s) = e^{sA}$ yields a one-parameter subgroup of $GL_n(\mathbb{C})$. \vskip{0.5em}

    In particular, the derivative is $\frac{d}{ds} \varphi_A(s) = A \varphi_A(s)$.
\end{boxed-text}

\begin{remark}
    (Reading Comprehension) Why is $\varphi_A(s) = e^{sA} \in GL_n(\mathbb{C})$ always true for any $A \in \mathbb{C}^{n \times n}$?
\end{remark}

And here's an even stronger observation.

\begin{theorem}{One-Parameter General Subgroups}
    Suppose $\varphi : \mathbb{R} \to GL_n(\mathbb{C})$ yields a one-parameter subgroup of $GL_n(\mathbb{C})$.  Then $\varphi$ looks like $\varphi: s \mapsto e^{As}$ for some $A \in \mathbb{C}^{n \times n}$.
\end{theorem}

\begin{proof}
    This is effectively a functional equation, with constraints:
    \[ (1) ~~ \varphi(0) = I_n ~~~~~~~ (2) ~~ \varphi(s + t) = \varphi(s) \varphi(t) \]
    Differentiating (2) and taking $s = 0$ yields $\varphi'(t) = \varphi'(0) \varphi(t)$.  Say $A = \varphi'(0)$.  The point is that $\varphi'(t) = A \varphi(t)$ is a linear first-order ODE, so it has a unique solution satisfying (1), and $\varphi: t \mapsto e^{tA}$ is such a solution. ~ $\blacksquare$
\end{proof}

\begin{remark}
    Recall the only solution to Cauchy's FE---provided $f$ is continuous!---is the linear equation.
\end{remark}

One can think of every one-parameter subgroup of $GL_n(\mathbb{C})$ as being parametrized by its ``tangent vector'' $\varphi'(0)$.  Once the ``initial conditions'' $\varphi(0) = I_n$ and $\varphi'(0) = A$ are set, the rest of the path is fixed!

\begin{center}
    \frame{\includegraphics[width=0.75\textwidth]{lecture-22/one-parameter-subgroups}}
\end{center}

\subsection{One-Parameter Subgroups of $SL_n(\mathbb{C})$ and $O_n(\mathbb{R})$}

More generally, for any subgroup $G \subseteq GL_n(\mathbb{C})$, the one-parameter subgroups of $G$ are precisely the images of maps $\varphi: \mathbb{R} \to G$ via $\varphi : s \mapsto e^{As}$ for some $A \in \mathbb{C}^{n \times n}$.  The only constraint on $A$ is that $e^{As} \in G$ for all $s \in \mathbb{R}$.

Let's investigate this constraint for various $G \subseteq GL_n(\mathbb{C})$.

\begin{theorem}{One-Parameter Special Subgroups}
    The one-parameter subgroups of $SL_n(\mathbb{C})$ are precisely the images of $\varphi: s \mapsto e^{sA}$ for some $A \in \mathbb{C}^{n \times n}$ satisfying $\mathrm{tr}(A) = 0$.
\end{theorem}

\begin{proof}
    We just need $e^{sA} \in SL_n(\mathbb{C})$ for all $s \in \mathbb{R}$, or:
    \[ \det e^{sA} = 1 ~ \implies ~ 1 = \det e^{sA} = e^{\mathrm{tr}(sA)} = e^{s \cdot \mathrm{tr}(A)}. \]
    And the above could only hold for all $s$ if $\mathrm{tr}(A) = 0$. ~ $\blacksquare$
\end{proof}

Similarly, we can talk about one-parameter subgroups of $O_n(\mathbb{R})$.

\begin{theorem}{One-Parameter Orthogonal Subgroups}
    The one-parameter subgroups of $O_n(\mathbb{R})$ are precisely the images of $\varphi: s \mapsto e^{sA}$ for some $A \in \mathbb{R}^{n \times n}$ satisfying $A^{\top} = -A$ (i.e. \textbf{skew-symmetric} matrices).
\end{theorem}

\begin{proof}
    We just need $e^{sA} \in O_n(\mathbb{R})$ for all $s \in \mathbb{R}$, or:
    \[ I_n = e^{sA}(e^{sA})^{\top} = e^{sA} e^{s A^{\top}}. \]
    Differentiating at $s = 0$ forces $A + A^{\top} = 0$, so $A^{\top} = -A$. ~ $\blacksquare$
\end{proof}

The one-parameter subgroups of $O_n(\mathbb{R})$ actually tell us about its dimension!

\begin{theorem}{Dimension of Orthogonal Subgroups}
    We have $\dim O_n(\mathbb{R}) = \binom{n}{2}$.
\end{theorem}

\begin{proof}
    It's equivalent to know the dimension of the set of one-parameter subgroups of $O_n(\mathbb{R})$.  But these biject with the set of skew-symmetric matrices, which is of dimension $\binom{n}{2}$. ~ $\blacksquare$
\end{proof}

\begin{remark}
    This is pretty hard to show without inspecting the one-parameter subgroups!
\end{remark}

Note that if $\varphi: s \mapsto e^{sA}$ describes a one-parameter subgroup of $G$, we need not have $A \in G$.  For example, $A = \begin{bsmallmatrix} 1 & 0 \\ 0 & -1 \end{bsmallmatrix}$ generates a one-parameter subgroup of $SL_2(\mathbb{C})$, but $\begin{bsmallmatrix} 1 & 0 \\ 0 & -1 \end{bsmallmatrix}$ is not in $SL_2(\mathbb{C})$.  We only need $A \in \mathbb{C}^{n \times n}$.

\begin{remark}
    (Reading Comprehension) Show that the one-parameter subgroups of $SU_2$ are precisely the images of $\varphi: s \mapsto e^{sA}$ for some $A \in \mathbb{C}^{n \times n}$ satisfying $A^{\dagger} = -A$ and $\mathrm{tr}(A) = 0$.  For example, $A = \begin{bsmallmatrix} i & -6 \\ 6 & -i \end{bsmallmatrix}$ works.

    If $A \in \mathbb{C}^{n \times n}$ does not satisfy this, then the path $\varphi(s) := e^{sA}$ will ``fly off'' the surface $S^3$.
\end{remark}

\subsection{Lie Algebras and Lie Brackets}

\textbf{Definition (Lie Algebra).} If $G$ is a Lie group, then its \textbf{Lie algebra} $\mathrm{Lie}(G)$ is the tangent space at its identity---in other words, the set of all possible $\varphi'(0)$ for one-parameter subgroups $\varphi: \mathbb{R} \to G$.

\textbf{Definition (Lie Bracket).} Every Lie algebra of matrices is equipped with the \textbf{Lie bracket} $[A, B] := AB - BA$.

\textbf{Example.} The Lie algebra of $SL_n(\mathbb{C})$ is $\{A \in \mathbb{C}^{n \times n} : \mathrm{tr}(A) = 0\}$, and that of $O_n(\mathbb{R})$ is $\{ A \in \mathbb{R}^{n \times n} : A^{\top} = -A \}$.  Meanwhile, the Lie algebra of $GL_n(\mathbb{C})$ is just all of $\mathbb{C}^{n \times n}$.

In the specific case of matrix groups, we have the following alternative definition:

\textbf{Definition (Lie Algebra).} If $G \subseteq GL_n(\mathbb{C})$ is a Lie group, then its \textbf{Lie algebra} is:
\[ \mathrm{Lie}(G) := \{ A \in \mathbb{C}^{n \times n} \mid  e^{sA} \in G \text{ for all } s \in \mathbb{R}\}. \]
Let's check some closure properties of Lie algebras and the Lie bracket.

\begin{theorem}{Closure of Addition}
    For any $A, B \in \mathrm{Lie}(G)$, we have $A + B \in \mathrm{Lie}(G)$.
\end{theorem}

\begin{proof}
    The fact that $A, B \in \mathrm{Lie}(G)$ is equivalent to the fact that $e^{sA}$ and $e^{sB}$ are in $G$ for all $s \in \mathbb{R}$.  Therefore, their product $e^{sA}e^{sB} = I_n + (A + B)s + O(s^2)$ is in $G$ for all $s \in \mathbb{R}$, so $A + B$ is a valid direction in the tangent space at the identity, meaning $A + B \in \mathrm{Lie}(G)$. ~ $\blacksquare$
\end{proof}

\begin{theorem}{Closure of Lie Bracket}
    For any $A, B \in \mathrm{Lie}(G)$, we have $[A, B] \in \mathrm{Lie}(G)$.
\end{theorem}

\begin{proof}
    The idea is to consider the following expression:
    \begin{align*}
        e^{As} e^{Bs} e^{-As} e^{-Bs} = \ & \left(1 + As + \frac{A^2s^2}{2} + O(s^3) \right) \left(1 + Bs + \frac{B^2s^2}{2} + O(s^3) \right)\left(\dots \text{$e^{-As}$}\dots \right)\left(\dots \text{$e^{-Bs}$}\dots \right) \\
        = \ & 1 + (AB - BA)s^2 + O(s^3) ~ \text{ after much simplification.}    \end{align*}
    The above is in $G$ for all $s \in \mathbb{R}$.  Thus, sending $s \to 0$, we see $e^{s^2(AB - BA)}$ is in $G$.  Therefore, ``locally'' speaking, $AB - BA$ is a valid direction in the tangent space at the identity, so $AB - BA = [A, B]$ is in $\mathrm{Lie}(G)$. ~ $\blacksquare$
\end{proof}


\begin{remark}
    It is not the case that $A, B \in \mathrm{Lie}(G)$ implies $AB \in \mathrm{Lie}(G)$.  (Try to find a counterexample!)
\end{remark}

The Lie Bracket also has a bunch of identities, which one can check if they really want to.

\begin{theorem}{Lie Bracket Identities}
    The following identities are true for any $A, B, C \in \mathrm{Lie}(G)$ and $c \in \mathbb{R}$.
    \[ (1) ~~ [A + B, C] = [A, C] + [B, C]  ~~~~~ (2) ~~ [cA, B] = c[A, B] ~~~~~ (3) ~~ [A, B] = -[B, A] \] \vspace{-0.8cm}
    \[ (4) ~~ [A, [B, C]] + [B, [C, A]] + [C, [A, B]] = 0 \] \vspace{-0.6cm}
\end{theorem}

\begin{remark}
    Identity (4) is known as the \emph{Jacobi Identity}, equivalent to $[A, [B, C]] = [[A, B], C] + [B, [A, C]]$.
    
    There's a neat analogy to the product rule $\frac{d}{dt} \left [ f(t) g(t) \right ]  = \left [ \frac{d}{dt}f(t) \right ] g(t) + f(t) \left [ \frac{d}{dt} g(t) \right ]$.
\end{remark}

Every Lie group has a corresponding Lie algebra, found via ``differentiation''.  However, just as ``integration'' always has a $+C$ term, not every Lie algebra has a unique Lie group; uniqueness is only guaranteed up to coverings.

\textbf{Example.} Both $\mathbb{R}$ and $S^1$ have isomorphic Lie algebras, and $\mathbb{R} / \mathbb{Z} \cong S^1$ indicates their mutual covering group.

We care about Lie algebras because they turn geometry and physics problems (i.e. Lie groups) into linear algebra problems (i.e. Lie algebras).

\end{document}